\chapter{جمع‌بندي و نتيجه‌گيري و پیشنهادات}\label{chap:5}
در این فصل، پس از بررسی جامع مفاهیم پایه‌ای، مدل‌سازی‌های ریاضی، الگوریتم‌های بهینه‌سازی و ساختار شبکه‌های عصبی عمیق، به جمع‌بندی دستاوردها و نتایج حاصل از این مستند می‌پردازیم. همچنین با تکیه بر دانش نظری به‌دست‌آمده در فصول پیشین، پیشنهادهایی کاربردی و مسیرهایی نوین برای توسعه این پژوهش و انجام کارهای آینده ارائه خواهد شد.

\section{جمع‌بندی و نتیجه‌گیری}\label{sec:5-1}
هدف اصلی این مستند، ایجاد یک پیوند علمی و منطقی میان مفاهیم بنیادین ریاضیات، الگوریتم‌های محاسباتی و فناوری‌های نوین یادگیری ماشین بود. روند این پژوهش نشان داد که چگونه مسائلی که در نگاه اول صرفاً تئوری به نظر می‌رسند، می‌توانند پایه و اساس فناوری‌های پیچیده امروزی باشند. نتایج و دستاوردهای اصلی این پژوهش را می‌توان در محورهای زیر خلاصه نمود:

\begin{itemize}
	\item \textbf{شناخت الگوهای پنهان در ساختارهای ریاضی:} در فصل نخست نشان داده شد که مسائلی نظیر تجزیه اعداد صحیح و توزیع اعداد اول که پایه و اساس سامانه‌های رمزنگاری و امنیت دیجیتال هستند، دارای رفتارهای هندسی معناداری (مانند مارپیچ اولام) می‌باشند. درک این الگوها پیش‌نیاز توسعه الگوریتم‌های هوشمند است.
	
	\item \textbf{اهمیت مدل‌سازی و بهینه‌سازی ریاضی:} در فصول دوم و سوم، چارچوب‌های تحلیلی و روش‌های بهینه‌سازی مورد ارزیابی قرار گرفتند. بررسی توابع زیان مختلف (مانند زیان هیوبر در برابر خطای میانگین مربعات) و الگوریتم‌هایی نظیر گرادیان کاهشی و آدام ثابت کرد که انتخاب درست معیار ارزیابی و روش بهینه‌سازی، نقشی حیاتی در مقاومت مدل در برابر داده‌های پرت و سرعت همگرایی آن دارد.
	
	\item \textbf{قدرت شبکه‌های عصبی در حل مسائل پیچیده:} در فصل چهارم با بررسی ساختار یاخته‌های عصبی مصنوعی و شبکه‌های چندلایه، نشان داده شد که این مدل‌ها به کمک توابع فعال‌ساز غیرخطی و الگوریتم پس‌انتشار خطا، توانایی فوق‌العاده‌ای در استخراج الگوهای پنهان داده‌ها دارند. 
	
	\item \textbf{هم‌گرایی امنیت و هوش مصنوعی:} در نهایت، به عنوان یکی از مهم‌ترین نتایج، مشخص شد که معماری‌های عمیق می‌توانند به عنوان ابزاری قدرتمند برای تامین امنیت فضای سایبری به کار گرفته شوند. مدل‌های یادگیری عمیق با تحلیل ترافیک شبکه توانایی بالایی در شناسایی حملات پیچیده نظیر منع سرویس توزیع‌شده و تشخیص نفوذ دارند و می‌توانند جایگزین قدرتمندی برای روش‌های سنتیِ مبتنی بر امضا باشند.
\end{itemize}

به طور خلاصه، این مستند یک چارچوب نظری منسجم فراهم کرد که نشان می‌دهد گذار از امنیت سنتیِ مبتنی بر پیچیدگی‌های محاسباتی (مانند تجزیه اعداد) به سمت امنیت هوشمندِ مبتنی بر داده (شبکه‌های عصبی عمیق)، یک مسیر علمی ضروری و اجتناب‌ناپذیر است.

\section{پیشنهادها برای کارهای آینده}\label{sec:5-2}
با توجه به مباحث نظری مطرح شده و پتانسیل بالای شبکه‌های عصبی عمیق، موارد زیر به عنوان پیشنهادهایی برای توسعه و ادامه این پژوهش در کارهای آینده ارائه می‌گردد:

\begin{enumerate}
	\item \textbf{پیاده‌سازی عملی مدل‌های تشخیص نفوذ:} پیشنهاد می‌شود در پژوهش‌های آتی، مفاهیم تئوری این مستند به صورت عملی پیاده‌سازی شوند. به عنوان مثال، استفاده از یک مجموعه دادگان\LTRfootnote{Dataset} استاندارد ترافیک شبکه و آموزش یک شبکه عصبی عمیق برای تشخیص بلادرنگ\LTRfootnote{Real-time} حملات سایبری و مقایسه دقت\LTRfootnote{Accuracy} آن با روش‌های کلاسیک می‌تواند بسیار ارزشمند باشد.
	
	\item \textbf{استفاده از یادگیری عمیق در تحلیل رمزنگاری\LTRfootnote{Cryptanalysis}:} با توجه به الگوهای بررسی شده در مارپیچ اولام، یکی از ایده‌های جذاب برای آینده، تغذیه این الگوهای بصری و هندسیِ اعداد اول به شبکه‌های عصبی (به‌ویژه شبکه‌های عصبی پیچشی\LTRfootnote{CNN}) است تا مشخص شود آیا هوش مصنوعی قادر به کشف روابط جدیدی برای تسریع فرآیند تجزیه اعداد صحیح می‌باشد یا خیر.
	
	\item \textbf{بررسی معماری‌های پیشرفته‌تر شبکه‌های عصبی:} در این مستند تمرکز اصلی بر شبکه‌های عصبی چندلایه  بود. پیشنهاد می‌شود برای داده‌های توالی‌دار مانند ترافیک شبکه که وابستگی زمانی دارند، از معماری‌های پیشرفته‌تری نظیر شبکه‌های عصبی بازگشتی\LTRfootnote{RNN} یا حافظه طولانی کوتاه‌مدت\LTRfootnote{LSTM} استفاده شود.
	
	\item \textbf{ترکیب الگوریتم‌های فراابتکاری با یادگیری عمیق:} بررسی ترکیب الگوریتم‌های بهینه‌سازی تکاملی با روش‌هایی مانند آدام برای فرار از کمینه‌های محلی\LTRfootnote{Local Minima} و بهبود فرآیند آموزش شبکه‌های عصبی، می‌تواند موضوع یک پژوهش مستقل و نوآورانه باشد.
\end{enumerate}
